% Valorant Competitive Analysis - Expanded Detailed Draft
\documentclass[11pt]{article}
% Packages
\usepackage[utf8]{inputenc}
\usepackage{graphicx}
\usepackage{booktabs}
\usepackage{hyperref}
\usepackage{caption}
\usepackage{float}
\usepackage{geometry}
\usepackage{amsmath}
\usepackage{siunitx}
\usepackage{longtable}
\usepackage{enumitem}
\geometry{margin=1in}

\title{A Comprehensive Analysis of Competitive Valorant Performance}
\author{Automated Analysis Pipeline}
\date{March 27, 2026}

\begin{document}
\maketitle

\begin{abstract}
This document provides a thorough, reproducible analysis of a large competitive Valorant snapshot dataset. The pipeline covers ingestion and cleaning of raw files, exploratory data analysis (EDA), unsupervised segmentation (clustering), and a supervised classification baseline to identify high-performing player snapshots. We include detailed data schema, transformation rules, figures, tables, model hyperparameters, evaluation metrics, and reproducibility instructions so readers can reproduce or extend the analysis.
\end{abstract}

\tableofcontents
\clearpage

\section{Introduction}
Valorant is a team-based first-person shooter where individual performance is influenced by agent choice, map, team composition, and regional meta. This analysis is intended to:
\begin{itemize}
  \item Document the raw data structure and cleaning decisions.
  \item Provide descriptive statistics and visualizations that summarize the dataset.
  \item Segment player snapshots using clustering to reveal typical play-style archetypes.
  \item Build a simple predictive model (Random Forest) to discriminate "high-rating" snapshots and report feature importance.
  \item Provide fully reproducible steps and artifacts.
\end{itemize}

\section{Data and Ingestion}
\subsection{Raw layout}
Files were provided under a multi-level directory tree rooted at \texttt{bronze/} organized by event and region. Each leaf contains a \texttt{data.csv} snapshot with consistent columns describing player-level statistics for a snapshot.

\paragraph{What we mean by "snapshot"}
In this paper a "snapshot" refers to a single per-leaf CSV file (the file named \texttt{data.csv}) that captures player-level statistics for a specific combination of event, region, map, agent and capture date. Each snapshot represents the dataset extracted at that time and place; when the pipeline merges snapshots it parses metadata from the file path (for example \texttt{event\_id}, \texttt{region}, \texttt{map}, \texttt{agent}, \texttt{snapshot\_date}) and appends them as columns so every row can be traced back to its originating snapshot.

\subsection{Merged artifact}
All per-snapshot CSV files were merged into a single streaming CSV to avoid out-of-memory concatenation. The primary merged artifact is \texttt{analysis/merged\_full.csv} (approximately 702k rows). A validated sample file is provided at \texttt{analysis/merged\_sample.csv} for quick inspection.

\subsection{Schema and column descriptions}
Table~\ref{tab:schema} summarizes the principal columns used throughout this analysis. Derived columns are noted.

\begin{longtable}{@{}p{0.28\textwidth}p{0.68\textwidth}@{}}
\caption{Selected data schema and descriptions.}\label{tab:schema}\\
\toprule
Column & Description \\
\midrule
\endfirsthead
	oprule
Column & Description \\
\midrule
\endhead
\bottomrule
\endfoot
  \texttt{player\_id} & Unique player identifier (string or numeric). \\
  \texttt{event\_id} & Tournament/event identifier (coarse grouping). \\
  \texttt{region} & Geographic region string (e.g., \textit{na}, \textit{eu}). \\
  \texttt{map} & Map name for the snapshot. \\
  \texttt{agent} & Agent name played in the snapshot (e.g., Jett, Sova). \\
  \texttt{snapshot\_date} & Snapshot capture date (ISO string where present). \\
  \texttt{rounds\_played} & Number of rounds played in the snapshot (integer). \\
  \texttt{kills}, \texttt{deaths}, \texttt{assists} & Raw counting stats for the snapshot. \\
  \texttt{kills\_per\_round} & Kills divided by rounds played (derived). \\
  \texttt{first\_kills}, \texttt{first\_deaths} & First kill/first death counts. \\
  \texttt{average\_combat\_score} (ACS) & Aggregated combat score used as a performance indicator. \\
  \texttt{average\_damage\_per\_round} (ADR) & Average damage per round. \\
  \texttt{headshot\_percentage} & Fraction (0--1) of kills that were headshots. \\
  \texttt{clutches\_played}, \texttt{clutches\_won}, \texttt{clutch\_success\_percentage} & Clutch-related metrics; clutch counts sometimes missing for casual snapshots. \\
  \texttt{rating} & Composite rating used as the target variable in this analysis. \\
  \texttt{org} & Organization identifier or name when available (mixed type in raw data). \\
\end{longtable}

Notes: percent-formatted fields in the raw CSV (e.g., ``25\%'') were converted to numeric fractions in preprocessing. Mixed-type columns (notably \texttt{org}) were normalized to strings where possible; when normalization failed the raw value was preserved to avoid data loss.

\section{Preprocessing and Cleaning}
This section enumerates the concrete cleaning rules applied by \texttt{analysis/merge\_and\_clean.py} and the streaming merger.

\subsection{Parsing rules}
\begin{itemize}[noitemsep]
  \item Percent strings: remove trailing ``\%'' and convert to float multiplied by 0.01 (e.g., ``25\%'' \textrightarrow 0.25).
  \item Clutch fields: strings like ``3/7'' were parsed into \texttt{clutches\_won}=3 and \texttt{clutches\_played}=7; a derived \texttt{clutch\_success\_percentage} = clutches\_won / clutches\_played computed when denominator $>0$. 
  \item Missing values: empty strings and explicit NA markers were converted to \texttt{NaN} using pandas' coercion rules.
  \item Types: numeric columns coerced with \texttt{pd.to\_numeric(errors='coerce')}; failures left as \texttt{NaN}.
  \item Metadata: event, region, map, agent and snapshot date were parsed from the file path when present and stored as separate columns.
\end{itemize}

\subsection{Parquet safety}
An attempt to write the merged dataset as Parquet initially failed due to mixed types in the \texttt{org} column. To ensure safe conversion, either cast \texttt{org} to string or inspect unique values and map numerics to string labels before writing Parquet. For reproducibility we left the canonical merged artifact as CSV (streamed append) and preserved the original raw types in a copy for auditing.

\section{Exploratory Data Analysis (EDA)}
This section summarizes key EDA outputs and references the figures generated by \texttt{analysis/eda.py}.

\subsection{Missingness and coverage}
Figure~\ref{fig:missing} shows the percent missing by column. Several optional fields (clutch fields, \texttt{org}) have substantial missingness, while core match stats (kills, deaths, rounds\_played) are near-complete.

\begin{figure}[H]
  \centering
  \includegraphics[width=0.9\textwidth]{../analysis/figures_full/missing_percent_by_column.png}
  \caption{Percent missing values by column.}
  \label{fig:missing}
\end{figure}

\subsection{Distribution of rating}
The empirical distribution of the measured \texttt{rating} is shown in Figure~\ref{fig:dist_rating}. The bulk of snapshots cluster around a mean near 1.0 with a right tail (high-performing outliers).

\begin{figure}[H]
  \centering
  \includegraphics[width=0.8\textwidth]{../analysis/figures_full/dist_rating.png}
  \caption{Distribution of player \texttt{rating} across the dataset.}
  \label{fig:dist_rating}
\end{figure}

\subsection{Correlations}
Figure~\ref{fig:corr} visualizes pairwise correlations for core numeric features. Expected strong correlations exist between ACS, ADR, and rating; kills and KPR also correlate strongly with rating.

\begin{figure}[H]
  \centering
  \includegraphics[width=0.9\textwidth]{../analysis/figures_full/correlation_matrix.png}
  \caption{Correlation matrix for core numeric features.}
  \label{fig:corr}
\end{figure}

\subsection{Top agents and organizations}
The dataset is dominated by certain agents and organizations. Table~\ref{tab:top-agents} and Table~\ref{tab:top-orgs} list the top-10 entries by snapshot frequency (counts computed from the merged CSV).

\begin{table}[H]
  \centering
  \caption{Top 10 agents by snapshot count}
  \label{tab:top-agents}
  \begin{tabular}{lr}
  \toprule
    Agent & Snapshots \\
    \midrule
    jett & 67{,}106 \\
    sova & 66{,}477 \\
    omen & 64{,}139 \\
    killjoy & 53{,}650 \\
    raze & 49{,}505 \\
    viper & 44{,}023 \\
    breach & 42{,}116 \\
    cypher & 40{,}676 \\
    skye & 33{,}457 \\
    astra & 28{,}961 \\
    \bottomrule
  \end{tabular}
\end{table}

\begin{table}[H]
  \centering
  \caption{Top 10 organizations by snapshot count}
  \label{tab:top-orgs}
  \begin{tabular}{lr}
    	oprule
    Organization & Snapshots \\
    \midrule
    18 & 10{,}307 \\
    17 & 10{,}238 \\
    19 & 9{,}985 \\
    16 & 9{,}538 \\
    20 & 9{,}419 \\
    21 & 8{,}796 \\
    15 & 8{,}520 \\
    22 & 8{,}252 \\
    24 & 7{,}617 \\
    14 & 7{,}563 \\
    \bottomrule
  \end{tabular}
\end{table}

\section{Unsupervised Analysis: Clustering}
We applied KMeans clustering to a standardized set of core numeric features to discover recurring player snapshot archetypes. The steps are:
\begin{enumerate}
  \item Select numeric features with low missingness (ACS, ADR, KPR, headshot\_percentage, kills\_per\_round, assists\_per\_round, rounds\_played).
  \item Impute remaining missing numeric values with median values computed from the training sample.
  \item Scale features with \texttt{StandardScaler} (zero mean, unit variance).
  \item Evaluate K=2..8 using average silhouette score to pick the best K (silhouette maximized at $k=4$ in our run).
  \item Fit final KMeans with \texttt{random\_state}=42 for reproducibility and save cluster centers.
\end{enumerate}

Figure~\ref{fig:kmeans-pca} shows a PCA projection colored by cluster for a random sample.

\begin{figure}[H]
  \centering
  \includegraphics[width=0.8\textwidth]{../analysis/advanced_outputs/figures/kmeans_pca.png}
  \caption{KMeans clusters (k=4) projected to two principal components.}
  \label{fig:kmeans-pca}
\end{figure}

Table~\ref{tab:cluster-summary} shows a concise summary of clusters (counts and selected means) computed from \texttt{analysis/advanced\_outputs/cluster\_summary.csv}.

\begin{table}[H]
  \centering
  \caption{Cluster summary (selected statistics)}
  \label{tab:cluster-summary}
  \begin{tabular}{lrrrrrr}
    	oprule
    Cluster & Count & mean(rating) & mean(ACS) & mean(ADR) & mean(KPR) & mean(HSP) \\
    \midrule
    0 & 45{,}510 & 0.705 & 152 & 101 & 0.510 & 0.242 \\
    1 & 15{,}102 & 1.038 & 209 & 136 & 0.732 & 0.259 \\
    2 & 37{,}224 & 1.193 & 241 & 156 & 0.855 & 0.281 \\
    3 & 2{,}164  & 1.087 & 222 & 143 & 0.778 & 0.261 \\
    \bottomrule
  \end{tabular}
\end{table}

\subsection{Cluster interpretation}
Based on the center values and feature distributions:
\begin{itemize}
  \item Cluster 0: larger population with modest ACS and rating (likely average/bench snapshots).
  \item Cluster 1: solid performers with above-average ACS and moderate KPR.
  \item Cluster 2: highest-performing snapshots (highest ACS, ADR, KPR and rating) — aggressive fraggers / top-performers.
  \item Cluster 3: small cluster with high ACS but specific characteristics (may represent unusual sample contexts).
\end{itemize}

% Inserted cluster agent composition table (top agents per cluster)
\begin{table}[ht]
\centering
\caption{Top 3 agents (by snapshot count) per cluster — full merged data}
\begin{tabular}{rl}
\toprule
Cluster & Top agents \
\midrule
0 & omen, sova, breach \
1 & sova, omen, jett \
2 & jett, raze, sova \
3 & sova, omen, jett \
\bottomrule
\end{tabular}
\label{tab:cluster_agents_full}
\end{table}


\section{Supervised Baseline: Predicting High Rating}
\subsection{Problem framing}
We framed a binary prediction problem: \texttt{high\_rating} = 1 if a snapshot's \texttt{rating} is in the top 25\% quantile, otherwise 0. This choice produces a balanced but slightly imbalanced classification target suitable for an initial baseline.

\subsection{Modeling pipeline}
\begin{itemize}
  \item Feature set: numeric performance indicators (ACS, ADR, KPR, headshot\_percentage, kills, assists, rounds\_played, first\_kills, first\_deaths).
  \item Train/test split: stratified 80/20 split on \texttt{high\_rating} label.
  \item Imputation: median for numeric features computed on the training set.
  \item Scaling: StandardScaler fitted on training set only.
  \item Classifier: RandomForestClassifier(n\_estimators=200, max\_depth=12, random\_state=42) as a strong, interpretable baseline.
\end{itemize}

\subsection{Evaluation metrics}
We report ROC AUC, precision, recall, and an approximate confusion matrix on the test fold. Receiver operating characteristic and feature importances are saved to \texttt{analysis/advanced\_outputs/}.

\subsection{Results}
The Random Forest baseline achieved an AUC of approximately 0.982 on the held-out test fold (see Figure~\ref{fig:roc}). Feature importance ranking shows \texttt{kills\_per\_round}, \texttt{average\_combat\_score}, and \texttt{average\_damage\_per\_round} as top predictors (Figure~\ref{fig:fi}).

\begin{figure}[H]
  \centering
  \includegraphics[width=0.65\textwidth]{../analysis/advanced_outputs/figures/roc_curve.png}
  \caption{ROC curve for the \texttt{high\_rating} classifier (AUC \(\approx 0.982\)).}
  \label{fig:roc}
\end{figure}

\begin{figure}[H]
  \centering
  \includegraphics[width=0.75\textwidth]{../analysis/advanced_outputs/figures/feature_importances.png}
  \caption{Top feature importances from the Random Forest classifier.}
  \label{fig:fi}
\end{figure}

Full classification reports, the trained model artifact (\texttt{rf\_high\_rating.joblib}), and exact hyperparameters are in \texttt{analysis/advanced\_outputs/advanced\_report.txt} and related files.

\section{Discussion and Domain Interpretation}
Key takeaways:
\begin{itemize}
  \item Rating distribution indicates a central mass around typical match-level performance with a meaningful right tail of high performers.
  \item Headshot percentage, ACS, and damage metrics are meaningful predictors and align with intuitive player impact measures.
  \item Clustering reveals plausible archetypes; however, caution is warranted because snapshots mix match contexts and sample sizes across events/regions.
\end{itemize}

\subsection{Limitations}
\begin{itemize}
  \item Snapshot-level analysis lacks match-level context (team composition, opponent strength, map-side economy) which can confound interpretation.
  \item Organizational identifiers in the raw data contained mixed types causing downstream type coercion issues; careful auditing required before merging for production Parquet.
  \item Model is a baseline and does not use time-series or within-match structure that could improve predictive power.
\end{itemize}

\section{Reproducibility and Artifacts}
All scripts, data artifacts, and figure outputs used to produce this paper are in the repository under \texttt{analysis/} and \texttt{paper/}. Critical files:
\begin{itemize}
  \item \texttt{analysis/merged\_full.csv} — merged dataset (CSV, 702k rows).
  \item \texttt{analysis/eda.py} — script to produce EDA figures in \texttt{analysis/figures\_full/}.
  \item \texttt{analysis/advanced.py} — clustering and modeling pipeline; saves artifacts to \texttt{analysis/advanced\_outputs/}.
  \item \texttt{paper/valorant\_analysis.tex} — this LaTeX source.
\end{itemize}

\subsection{Environment}
The analysis was executed in a local Python environment. Minimal package list (see \texttt{analysis/requirements.txt}):
\begin{itemize}
  \item pandas
  \item numpy
  \item scikit-learn
  \item matplotlib
  \item seaborn
  \item joblib
\end{itemize}

To build the PDF locally run:
\begin{verbatim}
cd paper
pdflatex valorant_analysis.tex
\end{verbatim}

If \texttt{pdflatex} is not available, use an online LaTeX service or install TeXLive.

\section{Appendix: Additional tables and figures}
The \texttt{analysis/} directory contains many auxiliary CSVs and images. Notable artifacts:
\begin{itemize}
  \item \texttt{analysis/top\_agents.csv} and \texttt{analysis/top\_orgs.csv} — raw counts used above.
  \item \texttt{analysis/advanced\_outputs/cluster\_summary.csv} — full cluster summaries (means, std, counts) used to create Table~\ref{tab:cluster-summary}.
  \item \texttt{analysis/advanced\_outputs/rf\_high\_rating.joblib} — trained Random Forest model object.
\end{itemize}

\section{Future Work}
Possible next steps to strengthen the analysis:
\begin{itemize}
  \item Add match-level linkage to contextualize snapshots (team mates, opponent ratings, map economy).
  \item Use longitudinal snapshots (if available) to model player trajectories with temporal models (RNNs or temporal point processes).
  \item Replace static imputation with model-based imputation for missing fields, and explore causal inference for agent effects using matched controls.
\end{itemize}

\section*{Acknowledgements}
This analysis was produced by an automated pipeline; reproducible code and artifacts are included in the repository for inspection.

\end{document}
